% ===== Джерела до розділу 2. Ключі з префіксом ch02: =====
% Позиції ядра (bib-core.tex) тут НЕ дублюються.

% --- Законодавство України ---
\bibitem{ch02:kk361} Кримінальний кодекс України: статті 361 (несанкціоноване втручання в роботу ЕОМ), 361-1 (створення з метою використання, поширення або збуту шкідливих програмних чи технічних засобів), 361-2, 362, 363, 363-1. URL: \url{https://zakon.rada.gov.ua/laws/show/2341-14}.

\bibitem{ch02:kpk} Кримінальний процесуальний кодекс України від 13.04.2012 № 4651-VI (щодо вилучення й дослідження електронних доказів). URL: \url{https://zakon.rada.gov.ua/laws/show/4651-17}.

\bibitem{ch02:zu-ek} Про електронні комунікації: Закон України від 16.12.2020 № 1089-IX. URL: \url{https://zakon.rada.gov.ua/laws/show/1089-20}.

\bibitem{ch02:kmu-943} Про затвердження Загальних вимог до кіберзахисту об'єктів критичної інфраструктури та Порядку їх застосування: постанова Кабінету Міністрів України. URL: \url{https://cip.gov.ua/}.

\bibitem{ch02:ndtzi-1-4} НД ТЗІ 1.4-001-2000. Типове положення про службу захисту інформації в автоматизованій системі. Київ: ДСТСЗІ СБ України, 2000.

% --- Право ЄС та міжнародне ---
\bibitem{ch02:eccc} Regulation (EU) 2021/887 establishing the European Cybersecurity Industrial, Technology and Research Competence Centre and the Network of National Coordination Centres. \emph{Official Journal of the EU}. L 202. 2021.

\bibitem{ch02:csa} Regulation (EU) 2019/881 on ENISA and on information and communications technology cybersecurity certification (Cybersecurity Act). \emph{Official Journal of the EU}. L 151. P.~15--69. URL: \url{https://eur-lex.europa.eu/eli/reg/2019/881/oj}.

\bibitem{ch02:budapest-ap} Second Additional Protocol to the Convention on Cybercrime on enhanced co-operation and disclosure of electronic evidence. CETS No.~224. Strasbourg: Council of Europe, 2022.

\bibitem{ch02:assoc} Угода про асоціацію між Україною та Європейським Союзом: ратифіковано Законом України від 16.09.2014 № 1678-VII. URL: \url{https://zakon.rada.gov.ua/laws/show/984_011}.

% --- Стандарти та бази знань ---
\bibitem{ch02:iso27032} ISO/IEC 27032:2023. Cybersecurity --- Guidelines for Internet security. Geneva: ISO, 2023. URL: \url{https://www.iso.org/standard/76070.html}.

\bibitem{ch02:cvss} FIRST. Common Vulnerability Scoring System Version 4.0: Specification Document. Forum of Incident Response and Security Teams, 2023. URL: \url{https://www.first.org/cvss/v4.0/}.

\bibitem{ch02:taxii} TAXII Version 2.1: Trusted Automated Exchange of Intelligence Information. OASIS Standard, 2021. URL: \url{https://docs.oasis-open.org/cti/taxii/v2.1/}.

\bibitem{ch02:misp} MISP: Open Source Threat Intelligence Platform \& Standards for Threat Information Sharing. URL: \url{https://www.misp-project.org/}.

\bibitem{ch02:cis} CIS Critical Security Controls, Version 8.1. Center for Internet Security, 2024. URL: \url{https://www.cisecurity.org/controls/v8-1}.

% --- Українські наукові джерела ---
\bibitem{ch02:polotai-ids} Полотай О. І. та ін. Інтеграція IDS/IPS у корпоративні мережі та перевірка playbook на кіберполігонах. \emph{Кібербезпека: освіта, наука, техніка}. 2025. URL: \url{https://csecurity.kubg.edu.ua/index.php/journal/article/view/889}.

\bibitem{ch02:shevchenko-class} Шевченко С. М., Жданова Ю. Д., Гаркушенко А. М. Аналіз підходів до класифікації сучасних загроз інформаційній безпеці. \emph{Кібербезпека: освіта, наука, техніка}. 2025. URL: \url{https://csecurity.kubg.edu.ua/index.php/journal/article/view/945}.
